\hypertarget{chapter-17-c14-functions-and-lambdas}{%
\section{CHAPTER 17: C14 FUNCTIONS AND
LAMBDAS}\label{chapter-17-c14-functions-and-lambdas}}

\hypertarget{c14-functions-lambdas}{%
\section{C++14 FUNCTIONS \& LAMBDAS}\label{c14-functions-lambdas}}

C++11 introduced lambdas, but C++14 made them ``First Class Citizens.''
They gained the ability to be generic and to handle move-only types,
making them indispensable for modern asynchronous and functional
programming.

\hypertarget{generic-lambdas}{%
\subsection{1. Generic Lambdas}\label{generic-lambdas}}

In C++11, lambda parameters required concrete types. C++14 allows
\passthrough{\lstinline!auto!} parameters, making the lambda's call
operator a template.

\hypertarget{the-internal-mechanics}{%
\subsubsection{1.1 The Internal
Mechanics}\label{the-internal-mechanics}}

When you write a generic lambda, the compiler generates a closure object
with a templated \passthrough{\lstinline!operator()!}.

\begin{lstlisting}[language={C++}]
auto sum = [](auto a, auto b) {
    return a + b;
};

// Effectively becomes:
struct __lambda_unique_name {
    template<typename T, typename U>
    auto operator()(T a, U b) const {
        return a + b;
    }
};
\end{lstlisting}

\hypertarget{polymorphic-behavior}{%
\subsubsection{1.2 Polymorphic Behavior}\label{polymorphic-behavior}}

Generic lambdas enable elegant, type-agnostic code without the
boilerplate of traditional templates.

\begin{lstlisting}[language={C++}]
auto printer = [](const auto& container) {
    for (const auto& item : container) {
        std::cout << item << " ";
    }
    std::cout << "\n";
};

std::vector<int> v = {1, 2, 3};
std::list<std::string> l = {"A", "B"};

printer(v); // Works for vector
printer(l); // Works for list
\end{lstlisting}

\hypertarget{lambda-init-capture-generalized-capture}{%
\subsection{2. Lambda Init-Capture (Generalized
Capture)}\label{lambda-init-capture-generalized-capture}}

This is arguably the most important lambda upgrade. It allows you to
create new variables in the capture clause, and more importantly, it
enables \textbf{capturing move-only types} like
\passthrough{\lstinline!std::unique\_ptr!}.

\hypertarget{moving-into-a-lambda}{%
\subsubsection{2.1 Moving into a Lambda}\label{moving-into-a-lambda}}

In C++11, you couldn't move a \passthrough{\lstinline!unique\_ptr!} into
a lambda without ugly workarounds. C++14 solves this.

\begin{lstlisting}[language={C++}]
auto data = std::make_unique<LargeBuffer>();

// Capture by move: 'p' is initialized by moving 'data'
auto task = [p = std::move(data)]() {
    p->process();
};

// 'data' is now null; 'p' lives inside the lambda object
\end{lstlisting}

\hypertarget{renaming-captures}{%
\subsubsection{2.2 Renaming Captures}\label{renaming-captures}}

You can also rename variables or capture the result of an expression.

\begin{lstlisting}[language={C++}]
int x = 10;
auto check = [val = x + 5](int input) {
    return input > val;
};
\end{lstlisting}

\hypertarget{return-type-deduction-decltypeauto}{%
\subsection{\texorpdfstring{3. Return Type Deduction \&
\texttt{decltype(auto)}}{3. Return Type Deduction \& decltype(auto)}}\label{return-type-deduction-decltypeauto}}

C++14 expanded return type deduction to all functions, not just lambdas.

\hypertarget{rules-for-auto-return-types}{%
\subsubsection{\texorpdfstring{3.1 Rules for \texttt{auto} Return
Types}{3.1 Rules for auto Return Types}}\label{rules-for-auto-return-types}}

The function body must be visible to the compiler at the call site. If
there are multiple \passthrough{\lstinline!return!} statements, they
must all deduce to the same type.

\begin{lstlisting}[language={C++}]
auto get_value(bool flag) {
    if (flag) return 42;    // Deduces int
    else      return 0;     // Deduces int
    // return 3.14;         // ERROR: inconsistent types (int vs double)
}
\end{lstlisting}

\hypertarget{the-decltypeauto-powerhouse}{%
\subsubsection{\texorpdfstring{3.2 The \texttt{decltype(auto)}
Powerhouse}{3.2 The decltype(auto) Powerhouse}}\label{the-decltypeauto-powerhouse}}

Standard \passthrough{\lstinline!auto!} return type deduction uses
template argument deduction rules, which means \textbf{references are
stripped (decayed)}. \passthrough{\lstinline!decltype(auto)!} preserves
the exact type, including references and const-qualifiers.

\begin{lstlisting}[language={C++}]
int global_val = 100;

int& get_ref() { return global_val; }

// Returns by value (int)
auto proxy1() { return get_ref(); }

// Returns by reference (int&) - Perfect Forwarding of Return Type
decltype(auto) proxy2() { return get_ref(); }

void test() {
    proxy1() = 200; // ERROR: modifying a temporary
    proxy2() = 200; // SUCCESS: modifies global_val
}
\end{lstlisting}

\textbf{Deep Dive:} Use \passthrough{\lstinline!decltype(auto)!}
primarily in wrapper functions or generic code where you want to pass
through the return type of another function exactly as it is, without
knowing whether it returns by value or reference.
